% Part: sets-functions-relations
% Chapter: sets
% Section: pairs-and-products

\documentclass[../../../include/open-logic-section]{subfiles}

\begin{document}

\olfileid{sfr}{set}{pai}
\olsection{序对、元组与笛卡尔积}

\begin{explain}
根据外延性原理，集合的元素是没有顺序的。因此，如果要表示顺序，我们使用\emph{序对} $\tuple{x, y}$。在无序对 $\{x, y\}$ 中，顺序无关紧要：$\{x, y\} = \{y, x\}$。但在序对中则不然：如果 $x \neq y$，那么 $\tuple{x, y} \neq \tuple{y, x}$。

我们该怎样在集合论中理解序对呢？关键在于，我们要保持以下性质：两个序对相同，当且仅当它们具有相同的第一分量和相同的第二分量，即：
\[
  \tuple{a, b}= \tuple{c, d}\text{ 当且仅当 }a = c \text{ 且 }b=d.
\]。
我们可以利用 Wiener--Kuratowski 定义在集合论中定义序对。
\end{explain}

\begin{defn}[序对]\ollabel{wienerkuratowski}
	$\tuple{a, b} = \{\{a\}, \{a, b\}\}$。
\end{defn}

\begin{prob}
	使用 \olref[sfr][set][pai]{wienerkuratowski}，证明 $\tuple{a,
	b}= \tuple{c, d}$ 当且仅当 $a = c$ 且 $b=d$。
\end{prob}

\begin{explain}
在确定了序对的定义之后，我们就可以用它来定义更多的集合。例如，有时我们也会需要多于两个对象的有序序列，比如\emph{三元组} $\tuple{x, y, z}$，\emph{四元组} $\tuple{x, y, z, u}$，等等。我们可以把三元组看作特殊的序对，其中第一个元素本身就是一个序对：$\tuple{x, y, z}$ 是 $\tuple{\tuple{x, y},z}$。四元组同理：$\tuple{x,y,z,u}$ 是 $\tuple{\tuple{\tuple{x,y},z},u}$，以此类推。一般地，我们会讨论\emph{有序 $n$ 元组} $\tuple{x_1, \dots, x_n}$。

由序对或其他有序 $n$ 元组构成的某些集合将会非常有用。
\end{explain}

\begin{defn}[笛卡尔积]
给定集合 $A$ 和 $B$，它们的\emph{笛卡尔积} $A \times B$ 定义为
\[
  A \times B = \Setabs{\tuple{x, y}}{x \in A \text{ 且 } y \in B}.
\]。
\end{defn}

\begin{ex}
如果 $A = \{0, 1\}$，并且 $B = \{1, a, b\}$，那么它们的笛卡尔积是
\[
A \times B = \{ \tuple{0, 1}, \tuple{0, a}, \tuple{0, b},
    \tuple{1, 1}, \tuple{1, a}, \tuple{1, b} \}.
\]。
\end{ex}

\begin{ex}
如果 $A$ 是一个集合，那么 $A$ 与自身的乘积 $A \times A$ 也记作~$A^2$。它是由\emph{所有}满足 $x, y \in A$ 的序对 $\tuple{x, y}$ 构成的集合。所有三元组 $\tuple{x, y, z}$ 构成的集合是 $A^3$，以此类推。我们可以给出一个递归定义：
\begin{align*}
  A^1 & = A\\
  A^{k+1} & = A^k \times A
\end{align*}。
\end{ex}

\begin{prob}
列出 $\{1, 2, 3\}^3$ 的所有元素。
\end{prob}

\begin{prop}\ollabel{cardnmprod}
如果 $A$ 有 $n$ 个元素且 $B$ 有 $m$ 个元素，那么 $A \times B$ 有 $n\cdot m$ 个元素。
\end{prop}

\begin{proof}
对于 $A$ 中的每个元素~$x$，都有 $m$ 个形如 $\tuple{x, y} \in A \times B$ 的元素。令 $B_x = \Setabs{\tuple{x, y}}{y
  \in B}$。因为只要 $x_1 \neq x_2$，就有 $\tuple{x_1, y} \neq
\tuple{x_2, y}$，所以 $B_{x_1} \cap B_{x_2} = \emptyset$。而如果 $A = \{x_1,
\dots, x_n\}$，那么 $A \times B = B_{x_1} \cup \dots \cup B_{x_n}$，因此它有 $n\cdot m$ 个元素。

为了直观理解，可以将 $A \times B$ 的元素排列成网格：
\[
\begin{array}{rcccc}
  B_{x_1} = & \{\tuple{x_1, y_1} & \tuple{x_1, y_2} & \dots & \tuple{x_1, y_m}\}\\
  B_{x_2} = & \{\tuple{x_2, y_1} & \tuple{x_2, y_2} & \dots & \tuple{x_2, y_m}\}\\
  \vdots & & \vdots\\
  B_{x_n} = & \{\tuple{x_n, y_1} & \tuple{x_n, y_2} & \dots & \tuple{x_n, y_m}\}
\end{array}
\]。
由于 $x_i$ 各不相同，$y_j$ 也各不相同，网格中的序对互不相同，并且总共有 $n\cdot m$ 个。
\end{proof}

\begin{prob}
通过对~$k$ 进行归纳，证明对于所有 $k \ge 1$，如果 $A$ 有 $n$ 个元素，那么 $A^k$ 有 $n^k$ 个元素。
\end{prob}

\begin{ex}
如果 $A$ 是一个集合，那么 $A$ 上的一个\emph{单词}是指 $A$ 中元素的任意序列。一个序列可以看作 $A$ 中元素的一个 $n$ 元组。例如，如果 $A = \{a, b, c\}$，那么序列 ``$bac$'' 可以看作是三元组~$\tuple{b, a, c}$。单词，即符号序列，在计算机科学中至关重要。按照惯例，我们将 $A$ 的元素视为长度为~$1$ 的序列，将 $\emptyset$ 视为长度为~$0$ 的序列。那么，$A$ 上\emph{所有}单词构成的集合就是
\[
A^* = \{\emptyset\} \cup A \cup A^2 \cup A^3 \cup \dots
\]。
\end{ex}

\end{document}